\begin{document}

% Project-Specific Settings
\newcommand{\thesubject}{Informatik}
\newcommand{\thetopic}{Informatik und Gesellschaft}
\newcommand{\theplace}{Winterthur}

\settoggle{oneauthor}{true}
\newcommand{\mainauthor}{Cyril Blum}
\newcommand{\mainauthormail}{cyril.blum@edu.zh.ch}


\lhead[ \leftmark ]{\thetopic}
\rhead[\thesubject]{\textcopyleft{} \href{mailto:\mainauthormail}{\color{black}{\mainauthor}}, \thesubject, \the\year}

\pagestyle{empty}
\newgeometry{margin=1cm}
\footnotesize

\begin{landscape}
	\part*{Hannah Fry - Hello World}
	\section{Themenübersicht}
	Untenan finden Sie eine Übersicht einiger (Unter-)Themen, die im Buch von Hannah Fry behandelt werden. Es handelt sich hierbei um eine nicht abschliessende Auflistung. Weitere Themen können in Absprache mit der Lehrperson vereinbart werden. Beachten Sie zudem, dass die Themen nicht zwingend Unter-Kapiteln im Buch entsprechen, sondern dass gewisse Themen über ein Kapitel \enquote{verteilt} immer wieder auftauchen.
	\label{sec:themen}

	% Please add the following required packages to your document preamble:
	% \usepackage{booktabs}
	% \usepackage{multirow}
	%\begin{table}[H]
	\begin{longtable}{@{}r|p{5.3cm}||p{10cm}|p{10cm}@{}}
		\toprule
		\textbf{Kapitel}              & \textbf{Fragen / Unterthema}               & \textbf{Auftrag}                                                                                                                                                                                                                                                                                                                                                            & \textbf{Recherche-Auftrag                                                                                                                                                                                                                                                      } \\ \midrule\midrule
		\multirow{3}{*}{\textbf{Macht}}        & Algorithmen: Definition \& Kategorien & \begin{tabular}[c]{@{}l@{}}Was ist ein \textbf{Algorithmus}? Wie wird er definiert?\\ Welche \textbf{vier Kategorien} von Algorithmen nennt das Buch?\end{tabular}                                                                                                                                                              & Finden Sie für jede im Buch erwähnte Algorithmus-Kategorie ein weiteres, eigenes Beispiel                                                                                                                                                                                        \\\cmidrule(l){2-4}
		                              & Ethische Folgen von Algorithmen       & Welche \textbf{negativen Folgen} werden im Zusammenhang mit dem Umgang mit Algorithmen genannt? In welchen \textbf{Lebensbereichen / Situationen} sollten Algorithmen nicht oder mit Vorsicht angewandt werden?                                                                                                                                                                               & Führen Sie eine eigene Recherche zu heiklen Algorithmen durch, präsentieren Sie Ihre Resultate der Klasse.                                                                                                                                                                       \\\cmidrule(l){2-4}
		                              & Eigene Einschätzung \& Fazit          & Fassen Sie die \textbf{Kernaussagen} des Kapitels zusammen. Was nehmen Sie vom Lesen des Kapitels mit? Welche anderen Aspekte sind Ihnen beim Lesen aufgefallen / geblieben? Was scheint Ihnen wichtig?                                                                                                                                                                              & Recherchieren Sie, inwiefern ChatGPT einem der erwähnten Algorithmen zugeordnet werden kann oder nicht. Haben wir es mit einer neuen Kategorie von Algorithmus zu tun?                                                                                                           \\\midrule\midrule
		\multirow{5}{*}{\textbf{Daten}}        & Wert von Daten                        & Welche Aspekte von Daten werden im Kapitel als \textbf{wertvoll} bezeichnet?                                                                                                                                                                                                                                                                                                         & Recherchieren Sie die 5 wertvollsten Firmen der Welt und Beschreiben Sie, inwiefern Daten für diese Firmen eine Rolle spielen.                                                                                                                                                   \\\cmidrule(l){2-4}
		                              & Datenschutz                           & Was sind \textbf{anonymisierte Daten}? Welche Informationen können gemäss Buch selbst aus anonymisierten Daten herausgelesen werden und wie?                                                                                                                                                                                                                                                 & Recherchieren Sie, was das Internet über Sie weiss. Suchen Sie insbesondere nach \enquote{Data Breaches} in der Schweiz. Präsentieren Sie die Ergebnisse der Klasse                                                                                                                    \\\cmidrule(l){2-4}
		                              & Mikromanipulation \& Psychologie      & Was ist \textbf{Mikro-Manipulation}? Wie werden wir im Alltag durch Daten (mikro-)manipuliert?                                                                                                                                                                                                                                                                                       & Finden Sie 5 Lebensbereiche, in denen Sie manipuliert werden.                                                                                                                                                                                                                    \\\cmidrule(l){2-4}
		                              & Politische Einflussnahme              & Was sagt das Buch dazu, inwiefern Daten \textbf{politische Prozesse beeinflussen} können?                                                                                                                                                                                                                                                                                            & Recherchieren und Präsentieren Sie die Rolle von Cambridge Analytica im ersten Trump-Wahlkampf (2016).                                                                                                                                                                           \\\cmidrule(l){2-4}
		                              & Bewertungen \& Bürger-Rating          & Fassen Sie Ihre Erkenntnisse sowie die Aussagen aus dem Buch zum Thema \textbf{Bewertungen} zusammen. Wie werden \textbf{Menschen durch Firmen und Regierungen bewertet}? Welche Interessen verfolgen diese dabei?                                                                                                                                                                            & Recherchieren Sie, in welchen weiteren Formen Menschen freiwillig oder unfreiwillig durch Firmen oder Regierungen bewertet werden (z.B. FICO-Score, Fichen-Skandal etc.).                                                                                                        \\\midrule\midrule
		\multirow{4}{*}{\textbf{Justiz}}       & Algorithmen in der Justiz             & Wie werden \textbf{Entscheidungsbäume} und \textit{Random Forests} durch das Buch beschrieben? Fassen Sie die Beispiele aus dem Buch zu diesem Thema zusammen.                                                                                                                                                                                                                       & Recherchieren Sie, was ein \textbf{Entscheidungsbaum} und was ein \textbf{Random Forest} ist. Fassen Sie zusammen, wie diese funktionieren und wie diese angewandt werden. Finden Sie ein weiteres Anwendungsbeispiel ausserhalb der Justiz.                                                       \\\cmidrule(l){2-4}
		                              & Falsche und richtige Urteile          & Stellen Sie die Problematik der \enquote{\textbf{falsch-positiven}}, \enquote{\textbf{richtig-positiven}}, \enquote{\textbf{falsch-negativen}} und \enquote{\textbf{richtig-negativen}} Vorhersagen / Urteile vor und beurteilen Sie das Thema in Bezug auf die Justiz. Erklären Sie ebenfalls die Begriffe \enquote{\textbf{Spezifizität}} und \enquote{\textbf{Sensitivität}}. Stellen Sie die Problematik der \enquote{\textbf{maschinellen Vorurteile}} klar.                    & Recherchieren Sie mindestens 2 Beispiele von falsch-positiven oder falsch-negativen Entscheidungen durch die AI. Beschreiben Sie, ob diese Fälle generelle Zweifel an der AI aufkommen lassen, oder ob es sich um Einzelfälle handelt.                                           \\\cmidrule(l){2-4}
		                              & Schwierige Entscheide                 & Illustrieren Sie das Problem der \enquote{schwierigen Entscheide} anhand der Denkweisen, die im Buch \enquote{\textbf{Schnelles Denken, langsames Denken}} von Daniel Kahneman erwähnt werden.                                                                                                                        & Finden Sie weitere Beispiele von schwierigen Aufgaben, die \enquote{kontra-intuitives} Denken zur Findung der Lösung erfordern. Mit \enquote{kontra-intuitivem Denken} ist langsames, analytisches Denken gemeint. Bonus: können solche Aufgaben durch heutige AI gut bewältigt werden?      \\\midrule\midrule
		\multirow{3}{*}{\textbf{Medizin}}      & Mustererkennung                       & Erklären Sie, wie \textbf{maschinelle Mustererkennung} in der heutigen Medizin angewandt wird, indem Sie die Lern- und Funktionsweise von \enquote{\textbf{Neuronalen Netzwerken}} erklären.                                                                                                                                                                                                         & Illustrieren Sie, wofür \textbf{Konvolutive Neuronale Netzwerke} (Convolutional Neural Networks, CNNs) auch ausserhalb der Medizin genutzt werden können                                                                                                                                  \\\cmidrule(l){2-4}
		                              & Krebsdiagnosen \& AI                  & Erklären Sie, wie KI zur besseren Entdeckung von Krebszellen beitragen kann und welche \textbf{Herausforderungen sich betreffend Über- und Unter-Diagnose} stellen. Erklären Sie ebenfalls die Begriffe \enquote{\textbf{sensitiv}} und \enquote{\textbf{sepezifisch}} (siehe Kapitel zu Justiz). Welche \textbf{Herausforderungen stellen sich beim Trainieren von KI} für medizinische Diagnosen betreffend Trainings-Daten? & Finden Sie beispiele von Schweizer Spitälern, die KI bereits heute zu Krebsdiagnosezwecken verwenden.                                                                                                                                                                            \\\cmidrule(l){2-4}
		                              & Medizin \& Datenschutz                & Fassen Sie die \textbf{Datenschutz-Problematiken} zusammen, welche im Buch erwähnt werden. Was ist an medizinischen Daten anders / sensibler als an anderen Daten? Inwiefern könnten Sie diese Regelungen persönlich betreffen?  Erklären Sie die Problematik anhand kostenpflichtiger \textbf{Gentests} (Herkunfts-Tests).                                                                   & Finden Sie ein Beispiel aus der Schweiz, in welchem die digitale Speicherung von medizinischen Personendaten zu Bedenken aus Sicht des Datenschutz führt.                                                                                                                        \\\midrule\midrule
		\multirow{2}{*}{\textbf{Autos}}        & Autos \& Algorithmen                  & Erklären Sie, wie der \textbf{Algorithmus von Bayes} funktioniert und wie dieser für selbstfahrende Autos verwendet wird. Erklären Sie ebenfalls, weshalb es so schwierig ist, einen Algorithmus für selbstfahrende Autos \enquote{manuell} zu programmieren.                                                                                                                              & Finden Sie weitere Beispiele, in welchen Bayes angewandt wird. Recherchieren Sie zudem, ob Bayes auch heute noch angewandt wird für selbstfahrende Autos.                                                                                                                        \\\cmidrule(l){2-4}
		                              & Autos \& Ethik                        & Erläutern Sie einige \textbf{ethische Problematiken}, die mit selbstfahrenden Autos verbunden sind. Erklären Sie insbesondere die Problematik schwieriger ethischer Entscheidungen im Falle eines sicheren, nahenden Unfalls (FahrerIn oder PassantIn schützen?) sowie die \textbf{Problematik abnehmender Fahrer-}Fähigkeiten.                                                               & Finden Sie mindestens 5 weitere Problematiken, die mit selbstfahrenden Autos verbunden sind. Sofern Sie Ihre Überlegungen schlüssig und überzeugend darlegen, dürfen Sie hier auch Ihre eigenen Überlegungen anstellen statt zu recherchieren.                                   \\\midrule\midrule
		\multirow{2}{*}{\textbf{Kriminalität}} & Kriminalität \& Algorithmen           & Erläutern Sie, wie Kriminalität mittels \textbf{PredPol} bekämpft wird und beschreiben Sie die Funktionsweise von PredPol. Welche Vorgänger-Systeme gab es, inwiefern haben diese versagt, und weshalb sind gerade räumliche \textbf{(geografische) Daten} besonders aufschlussreich?                                                                                                         & Recherchieren Sie, welche Polizei-IT-Systeme in Europa und in der Schweiz angewandt werden, um nach Kriminellen zu fahnden.                                                                                                                                                      \\\cmidrule(l){2-4}
		                              & Kriminalität \& Falschdiagnosen       & Erläutern Sie detailliert die Problematik der \textit{\textbf{doubles}} und fassen Sie die \textbf{Gesichtserkennungs-Algorithmen} zusammen, welche im Buch erwähnt werden.                                                                                                                                                                                                                            & Recherchieren Sie, wie sich der Social Credit Score / Sesame Score in China seit Veröffentlichung des Buches weiterentwickelt hat und präsentieren Sie die Resultate ausführlich (min. 3-4 Minuten).                                                                             \\\midrule\midrule
		\multirow{2}{*}{\textbf{Kunst}}        & Fake Kunst \& Manipulation            & Erläutern Sie das Thema \textbf{\enquote{\textit{social proof}}} anhand der im Buch aufgeführten Beispiele für Streaming-Dienste.                                                                                                                                                                                                                                                          & Finden Sie 2 weitere Beispiele für social proof mit Bezug zur Informatik. Dabei kann es sich auch um Beispiele aus der Finanzwelt handeln (z.B. Bitcoin \& Social Proof, Game Stop Hype \& Social Proof etc.).                                                                   \\\cmidrule(l){2-4}
		                              & Kunst \& Qualität                     & Erläutern Sie die Problematik der \textbf{Messung von Qualität in der Kunst} anhand der Beispiele aus dem Buch.                                                                                                                                                                                                                                                                      & Präsentieren Sie mindestens 3 weitere Beispiele aus der Kunst und erstellen Sie, wenn möglich, eigene Beispiele (z.B. Suno, Meta Imagine etc.). Sind die Beispiele originell / von hoher Qualität?                                                                               \\
	\end{longtable}

	%\end{table}
\end{landscape}


\restoregeometry
\pagestyle{fancy}
\newpage
\section{Auftrag}
Lesen Sie im Buch \enquote{Hello World} von Hannah Fry das Kapitel \enquote{Einleitung} und ein Kapitel Ihrer Wahl (siehe Themenübersicht in Abschnitt \ref{sec:themen}).

Selbstverständlich dürfen Sie auch weitere Kapitel aus dem Buch lesen. Machen Sie sich beim Lesen des Buchs Notizen zu den Kapiteln (worum geht es? Was ist Ihnen geblieben? Was verstehen Sie / nicht? etc.).

Im Anschluss an die Lektüre erstellen Sie in Gruppen eine \textbf{Präsentation zu Ihrem Kapitel} (siehe Abschnitt \ref{sec:themen}). Die Gruppengrössen variieren je nach Kapitel: Tragen Sie sich für ein Thema Ihrer Wahl ein (\href{https://eduzh-my.sharepoint.com/:x:/g/personal/cyril_wendl_edu_zh_ch/EScHfDarDeVHhAJB3WPCghsBYa63XpdaEkcrZk5dOiQiVA?e=kHy0bM}{Link}). Falls es zu viele InteressentInnen für ein Kapitel gibt, sprechen Sie sich bitte untereinander ab, oder kontaktieren Sie die Lehrperson, um eine mögliche Aufteilung des Kapitels in mehrere Unter-Kapitel zu besprechen.

\textbf{Anforderungen der Präsentation:}
\begin{itemize}
	\item Dauer: ca. \textbf{20 bis maximal 30 Minuten} (ohne Diskussion), die Redezeit der Gruppenmitglieder sollte dabei ungefähr gleich verteilt sein.
	\item Inhalt: Eine didaktisch hochwertige Zusammenfassung des von Ihnen behandelten Kapitels aus dem Buch von Hannah Fry sowie der Recherche-Aufträge (siehe \cref{sec:themen} und \cref{sec:materialien} für Recherche-Materialien). Die Unterthemen des Kapitels sowie die in der Tabelle \textbf{fett markierten} Aspekte sollten alle behandelt werden, inklusive Recherche-Aufträge.
	\item Die detaillierten Bewertungskriterien finden Sie im Bewertungsraster auf den folgenden Seiten.
\end{itemize}

Falls Sie wollen, dürfen sie gerne auch ausgefallenere und kreative Präsentationsformen wählen (z.B. Rollenspiel, Interview, Podcast-Format etc.). Sprechen Sie dies jedoch bitte vorher mit der Lehrperson ab. Falls Sie mögen, können Sie im Anschluss an die Präsentation eine kurze Frage- und Diskussionsrunde (max. 5 Minuten) mit der Klasse durchführen.

\newpage
\restoregeometry
\pagestyle{empty}
\newgeometry{margin=0.6cm}

\section*{Bewertungsraster Präsentation}

\footnotesize
\vspace{0.3cm}

\noindent
\newcommand{\headFld}[2]{\TextField[name=#1, width=#2, height=.35cm, borderwidth=0]{} }
\textbf{Gruppe:} \headFld{gruppe}{5.4cm} \quad \textbf{Datum:} \headFld{datum}{3.8cm} \quad \textbf{Kapitel:} \headFld{kapitel}{5.4cm}

\vspace{0.3cm}
\newcommand{\commentFld}[1]{\TextField[value={Kommentar:}, name=#1, width=15.5cm, height=1.9cm, borderwidth=0, charsize=5pt, multiline=true, maxlen=1000]{}}
\newcommand{\ptFld}[1]{\TextField[name=#1, width=1cm, height=.3cm, borderwidth=0]{}}

\setlength{\tabcolsep}{4pt}
\begin{longtable}{@{}p{2.4cm}p{13cm}cc@{}}
	\toprule
	\textbf{Kriterium} & \textbf{Beschreibung} & \textbf{Max.} & \textbf{Pkt.} \\
	\midrule\midrule
	
	\textbf{Buchinhalt:\newline Verständnis} & Kapitel aus Hannah Frys Buch \textbf{verstanden}, \textbf{korrekt wiedergegeben}, \textbf{Kernaussagen} klar & 3 & \ptFld{pts_buchinhalt} \\
	\cmidrule(l){2-4}
	& \multicolumn{3}{p{15.5cm}}{\commentFld{kommentarBuchinhalt}} \\
	\midrule
	
	\textbf{Recherche:\newline Qualität} & \textbf{Eigenständige}, \textbf{aktuelle}, \textbf{relevante} Recherche; über Buch hinausgehend, Bezug zu Kapitel klar & 3 & \ptFld{pts_recherche} \\
	\cmidrule(l){2-4}
	& \multicolumn{3}{p{15.5cm}}{\commentFld{kommentarRecherche}} \\
	\midrule
	
	\textbf{Aufbau \&\newline Struktur} & \textbf{Klare Gliederung}, \textbf{logischer Aufbau}, roter Faden erkennbar & 2 & \ptFld{pts_aufbau} \\
	\cmidrule(l){2-4}
	& \multicolumn{3}{p{15.5cm}}{\commentFld{kommentarAufbau}} \\
	\midrule
	
	\textbf{Didaktik} & \textbf{Verständlich} erklärt, \textbf{anschauliche Beispiele}, Komplexität reduziert & 3 & \ptFld{pts_didaktik} \\
	\cmidrule(l){2-4}
	& \multicolumn{3}{p{15.5cm}}{\commentFld{kommentarDidaktik}} \\
	\midrule
	
	\textbf{Fachsprache} & Fachbegriffe \textbf{korrekt verwendet} und \textbf{erklärt} & 2 & \ptFld{pts_fachsprache} \\
	\cmidrule(l){2-4}
	& \multicolumn{3}{p{15.5cm}}{\commentFld{kommentarFachsprache}} \\
	\midrule
	
	\textbf{Medien \&\newline Visualisierung} & \textbf{Eigenständig erstellt}, \textbf{übersichtlich}, \textbf{unterstützt Inhalt}; Abzug: kopiert (-3) & 3 & \ptFld{pts_medien} \\
	\cmidrule(l){2-4}
	& \multicolumn{3}{p{15.5cm}}{\commentFld{kommentarMedien}} \\
	\midrule
	
	\textbf{Auftreten \&\newline Vortrag} & \textbf{Frei} gesprochen, \textbf{sicher}, \textbf{engagiert}, angemessenes Tempo & 2 & \ptFld{pts_auftreten} \\
	\cmidrule(l){2-4}
	& \multicolumn{3}{p{15.5cm}}{\commentFld{kommentarAuftreten}} \\
	\midrule
	
	\textbf{Zeit-\newline Management} & \textbf{20--30 Minuten}, \textbf{gleichmässige Redeverteilung}; Abzug: zu kurz/lang (-2) & 0 & \ptFld{pts_zeit}	 \\
	\cmidrule(l){2-4}
	& \multicolumn{3}{p{15.5cm}}{\commentFld{kommentarZeit}} \\
	\midrule\midrule
	
	\multicolumn{3}{r}{\textbf{Gesamtpunktzahl (max. 20):}} & \ptFld{total} \\
	\bottomrule
\end{longtable}

\newpage
\restoregeometry
\pagestyle{fancy}

\section*{Anhang: Materialien zur Recherche}\label{sec:materialien}

Die folgenden Materialien können für die Recherche-Aufträge verwendet werden. Sie behandeln aktuelle Themen im Bereich künstliche Intelligenz und Digitalisierung und eignen sich zur Ergänzung der Buchkapitel. Alle Materialien stehen hier zur Verfügung:

\textbf{PDFs:} \href{https://cvw.synology.me:5001/d/s/10u2MOpsQuAiafpqNxks9URnYdmNhpO7/lmoTdm0b0TOQs0N1U4QiMzvCNmW0kESX-r7HgB4Js7gw}{\faFilePdf{}  NZZ-Artikel für Recherche-Aufträge} (\textbf{Passwort:} \texttt{NZZ})

\noindent
\textbf{Videos:} \href{https://cvw.synology.me:5001/d/s/11QxlSoKIKc3fd1xPi7rD4o9xUfriMVH/Zu7mRnEeIx1XPU-wvav24IlXmIu5I9D6-CLngzwnQ6gs}{\faPlayCircle{} Videosammlung Informatik} (\textbf{Passwort:} \texttt{FilmeInformatik})

\vspace{0.4cm}

\footnotesize
\begin{longtable}{@{}lp{14cm}@{}}
	\toprule
	\textbf{Kapitel} & \textbf{Passende Materialien (PDFs \& Videos)} \\
	\midrule\midrule
	
	\textbf{Macht} & 
	\begin{itemize}[leftmargin=*,noitemsep,topsep=0pt]
		\item[\faFilePdf{}] \textit{Die entscheidende Frage für 2026 lautet: Welche KI wollen wir?} (12. Januar 2026) \newline
		$\rightarrow$ KI-Entwicklung, gesellschaftliche Kontrolle, ethische Fragen
		\item[\faFilePdf{}] \textit{Mein Boss, die KI} (15. Februar 2026) \newline
		$\rightarrow$ KI-Agenten, Gig-Economy, Bruecke zwischen digitaler und physischer Welt
		\item[\faPlayCircle{}] \textit{ChatGPT und Co. -- Wo künstliche Intelligenz verblüfft und wo sie scheitert} (29. Juni 2023) \href{https://www.srf.ch/play/tv/einstein/video/chatgpt-und-co----wo-kuenstliche-intelligenz-verbluefft-und-wo-sie-scheitert?urn=urn:srf:video:3c267144-f388-4474-83fc-a8d09676c94a}{\faLink}
		\item[\faPlayCircle{}] \textit{Was kann KI wirklich?} (21. November 2024) \href{https://www.srf.ch/play/tv/einstein/video/was-kann-ki-wirklich?urn=urn:srf:video:9c65b2f1-abaa-4583-9ddf-17e1a7743c2d}{\faLink}
	\end{itemize} \\
	\midrule
	
	\textbf{Daten} & 
	\begin{itemize}[leftmargin=*,noitemsep,topsep=0pt]
		\item[\faFilePdf{}] \textit{Tech-Firmen müssen KI-Tools sicherer machen} (23. Mai 2024) \newline
		$\rightarrow$ Datenschutz, Trainingsdaten, illegale Inhalte, Verantwortung
		\item[\faFilePdf{}] \textit{Telegram ist für alles offen} (27. Juli 2024) \newline
		$\rightarrow$ Datenschutz, Verschlüsselung, Überwachung vs. Privatsphäre
		\item[\faFilePdf{}] \textit{Die entscheidende Frage für 2026 lautet: Welche KI wollen wir?} (12. Januar 2026) \newline
		$\rightarrow$ Datenverarbeitung, Datensouveränität, Privatsphäre
		\item[\faPlayCircle{}] \textit{Im Bann von TikTok, Instagram und Co.} (11. August 2024) \href{https://www.srf.ch/play/tv/kassensturz/video/im-bann-von-tiktok-instagram-und-co-?urn=urn:srf:video:3cdc375c-5f1d-4029-bda4-9690c27c9c0e}{\faLink}
		\item[\faPlayCircle{}] \textit{Hype um TikTok -- Lieben oder löschen?} (24. August 2024) \href{https://www.srf.ch/play/tv/dok/video/hype-um-tiktok---lieben-oder-loeschen?urn=urn:srf:video:6b2838d3-7f05-4dc3-897d-dfc68b7430c7}{\faLink}
		\item[\faPlayCircle{}] \textit{Aufwachsen ohne Smartphone -- Ist eine Kindheit ohne TikTok und Co. möglich?} (21. August 2025) \href{https://www.srf.ch/play/tv/nzz-format/video/aufwachsen-ohne-smartphone---ist-eine-kindheit-ohne-tiktok-und-co--moeglich?urn=urn:srf:video:0fb13fb2-cc18-448f-9a5f-2ff411c16877}{\faLink}
	\end{itemize} \\
	\midrule
	
	\textbf{Justiz} & 
	\begin{itemize}[leftmargin=*,noitemsep,topsep=0pt]
		\item[\faFilePdf{}] \textit{Tech-Firmen müssen KI-Tools sicherer machen} (23. Mai 2024) \newline
		$\rightarrow$ Falsch-positive/negative Entscheidungen, maschinelle Vorurteile
		\item[\faPlayCircle{}] \textit{Dein Gesicht gehört uns -- Die Gefahr automatisierter Gesichtserkennung} (29. Juni 2025) \href{https://www.srf.ch/play/tv/dok/video/dein-gesicht-gehoert-uns---die-gefahr-automatischer-gesichtserkennung?urn=urn:srf:video:6596261a-5936-4d5f-9c6c-51c985e91534}{\faLink}
	\end{itemize} \\
	\midrule
	
	\textbf{Medizin} & 
	\begin{itemize}[leftmargin=*,noitemsep,topsep=0pt]
		\item \textit{Hinweis:} Für dieses Kapitel sind keine spezifischen Artikel vorhanden. Eigenständige Recherche erforderlich (z.B. zu KI in Schweizer Spitälern, medizinische Bilderkennung).
		\item[\faPlayCircle{}] \textit{KI in der Medizin -- Sind Algorithmen bald die besseren Ärzte?} (16. November 2023) \href{https://www.srf.ch/play/tv/nzz-format/video/ki-in-der-medizin---sind-algorithmen-bald-die-besseren-aerzte?urn=urn:srf:video:0d2820b2-175f-4832-b51e-8b4008963a8a}{\faLink}
	\end{itemize} \\
	\midrule
	
	\textbf{Autos} & 
	\begin{itemize}[leftmargin=*,noitemsep,topsep=0pt]
		\item \textit{Hinweis:} Für dieses Kapitel sind keine spezifischen Artikel vorhanden. Eigenständige Recherche zu selbstfahrenden Autos und ethischen Dilemmata erforderlich.
		\item[\faPlayCircle{}] \textit{Autonomes Fahrzeug -- Mobilität der Zukunft oder Wunschdenken?} (11. September 2023) \href{https://www.srf.ch/play/tv/einstein/video/autonomes-fahrzeug-mobilitaet-der-zukunft-oder-wunschdenken?urn=urn:srf:video:e107b861-fb6b-48e0-9471-4f55de35498c}{\faLink}
	\end{itemize} \\
	\midrule
	
	\textbf{Kriminalität} & 
	\begin{itemize}[leftmargin=*,noitemsep,topsep=0pt]
		\item[\faFilePdf{}] \textit{Tech-Firmen müssen KI-Tools sicherer machen} (23. Mai 2024) \newline
		$\rightarrow$ Deepfakes, Gesichtserkennung, Missbrauch von KI-Tools
		\item[\faFilePdf{}] \textit{Telegram ist für alles offen} (27. Juli 2024) \newline
		$\rightarrow$ Kriminalität in Messaging-Apps, Überwachung, Strafverfolgung
		\item[\faPlayCircle{}] \textit{Dein Gesicht gehört uns -- Die Gefahr automatisierter Gesichtserkennung} (29. Juni 2025) \href{https://www.srf.ch/play/tv/dok/video/dein-gesicht-gehoert-uns---die-gefahr-automatischer-gesichtserkennung?urn=urn:srf:video:6596261a-5936-4d5f-9c6c-51c985e91534}{\faLink}
		\item[\faPlayCircle{}] \textit{Mit Roboter und Drohnen: Die mögliche Zukunft der Einsatzkräfte} (25. Mai 2022) \href{https://www.srf.ch/play/tv/einstein/video/mit-roboter-und-drohnen-die-moegliche-zukunft-der-einsatzkraefte?urn=urn:srf:video:4b1b9ca9-32f2-44f3-a4fd-5ee02b59cbfa}{\faLink}
	\end{itemize} \\
	\midrule
	
	\textbf{Kunst} & 
	\begin{itemize}[leftmargin=*,noitemsep,topsep=0pt]
		\item[\faFilePdf{}] \textit{Der Prompt schult das kreative Schreiben} (26. Januar 2026) \newline
		$\rightarrow$ KI-generierte Kunst, Qualitätsfrage, neue Kunstformen
		\item[\faSafari] Webseiten: Suno, Meta Imagine, DALL-E etc. \newline
		$\rightarrow$ KI-Kunstgeneratoren, Urheberrecht, Originalität
		\item[\faPlayCircle{}] \textit{Echt oder KI -- Ist unsere Stimme auswechselbar?} (8. Mai 2025) \href{https://www.srf.ch/play/tv/einstein/video/echt-oder-ki-ist-unsere-stimme-auswechselbar?urn=urn:srf:video:e7fc0175-a736-4e9e-8e98-f27f231175c4}{\faLink}
	\end{itemize} \\
	
	\bottomrule
\end{longtable}

\normalsize
\newpage
\subsection*{Hinweise zur Verwendung}

\begin{itemize}
	\item Die Artikel dienen als \textbf{Ausgangspunkt} für die Recherche-Aufträge. Sie sollten durch eigene, aktuelle Recherchen ergänzt werden.
	\item \textbf{Mindestquellen:} Mindestens \textbf{2 Quellen}, davon \textbf{mind. 1 Quelle jünger als 3 Jahre} und \textbf{mind. 1 Quelle mit Schweiz-Bezug}.
	\item \textbf{Eigenleistung:} Mindestens \textbf{eine selbst erstellte Grafik/Diagramm} auf Basis eigener Datenrecherche (keine kopierten Visuals/Screenshots).
	\item Achten Sie darauf, die Artikel \textbf{kritisch} zu lesen und verschiedene Perspektiven einzubeziehen.
	\item \textbf{Schweiz-Bezug:} Versuchen Sie, wo möglich, Verbindungen zur Schweizer Situation herzustellen (z.B. Datenschutzgesetze, Schweizer Firmen, nationale Diskussionen).
	\item Bei der Verwendung von Artikelinhalten im Fact Sheet oder Vortrag sind die \textbf{Quellen korrekt anzugeben}.
	\item Sie dürfen in Ihrer Präsentation einen substantiellen Anteil der Zeit auf die \textbf{Recherche-Ergebnisse} verwenden, um das Kapitel aus dem Buch zu ergänzen und zu vertiefen (bis zu 50\% der Präsentationszeit).
\end{itemize}

\end{document}